\scanpage{294}
\mspage{244}

La construction de ce foncteur $i_*$ devrait être essentiellement triviale, avec ce dont nous disposons. Nous allons \deleted{\unsure{le}} définir d'abord, \deleted{sur la sous-catégorie} pour \uline{$S$ fixé}, un foncteur canonique
\[
 Ho^{\unsure{lc|le}}_{W\deleted{\unsure{,}}}(S)\times Ho^{lc}_{W}(S)
 \longrightarrow Ho_{W,\deleted{\unsure{S'}}S}(\Phi) \tag*{$(*)$}
\]
où le deuxième membre est défini, par analogie avec les $Ho_W(I)$, par
\[
 Ho_{W,S}(I)\overset{\mathrm{def}}{=}
 \bigl(M_{/S}(I)\bigr)\bigl(W_S(I)\bigr)^{-1}
\]
-- i.e. c'est la $=$ définition, sauf que $(M,W)$ est remplacé par $(M_{/S},W_S)$. Quant à
\diagram{\begin{tikzcd}[ampersand replacement=\&,row sep=tiny,column sep=small]
 \&1\arrow[dr,"{\deleted{\unsure{\nwarrow}}}"{near start}]\&\\ \Phi={} \&\&3,\\ \&2\arrow[ur,"{\deleted{\unsure{\swarrow}}}"{near start}]\&
\end{tikzcd}}
on le considère comme une sous-catégorie fermée de $\Psi^{\sim}$, qui s'en déduit en lui rajoutant un objet initial $0$ (tout comme $\Psi^{\sim}$ se déduit de $\Psi$ en lui rajoutant un objet final $3$). D'autre part, on a un foncteur tautologique
\[
 Ho_{W,S}(\Phi)\longrightarrow Ho_W(\deleted{\unsure{\Phi}}\unsure{\Psi^{\sim}|\widetilde\Psi}),
 \tag*{$(*')$}
\]
qui sur les objets consiste à compléter un diagramme
\diagram{\begin{tikzcd}[ampersand replacement=\&,row sep=tiny,column sep=small]
 X_1\&\\ \&Y\arrow[ul]\arrow[dl]\\ X_2\&
\end{tikzcd}}
de $\deleted{\unsure{Cat}}M_{/S}$ en le diagramme carré
\diagram{\begin{tikzcd}[ampersand replacement=\&,row sep=tiny,column sep=small]
 \&X_1\arrow[dl]\&\\ S\&\&Y\arrow[ul]\arrow[dl]\\ \&X_2\arrow[ul]\&
\end{tikzcd}}
de $M$.

% archive sheet number: 294

